\begin{frame}{Об'єкт, предмет та мета дослідження}
    \begin{itemize}
        \item \textbf{Об'єкт дослідження.}
            Інтерактивна багатоваріантна колоризація зображень.

        \item \textbf{Предмет дослідження.}
            Методи глибокого навчання для колоризації зображень,
            зокрема згорткові нейронні мережі (CNN), генеративно-змагальні мережі (GAN),
            трансформери (Transformer), дифузійні моделі (Diffusion Models)
            та селективні моделі простору станів (SSM),
            а також методи та метрики для їх порівняльного аналізу.

        \item \textbf{Мета дослідження.}
            Розробити програмне забезпечення інтерактивної багатоваріантної колоризації зображень та модель колоризації на базі SSM,
            після чого провести комплексне порівняння цієї моделі з CNN, GAN, Transformer та Diffusion
            як за об'єктивними метриками, так і за результатами користувацького дослідження.
    \end{itemize}
\end{frame}

% \begin{frame}{Об'єкт, предмет та мета дослідження}
%     \begin{itemize}
%         \item \textbf{Об'єкт дослідження.}
%             % Інтерактивна багатоваріантна колоризація зображення.

%         \item \textbf{Предмет дослідження.}
%             % Актуальні методи колоризації зображення.
    	
%         \item \textbf{Мета дослідження.}
%             % Розробити програмне забезпечення колоризації зображень різними моделями.
%             % Розробити модель колоризації зображень на базі селективної моделі простору станів (англ. Selective State Space Model, SSM).
%             % Провести порівняння власної моделі з актуальними методами колоризації, оцінючи як об'єктивно через метрики, так і суб'єктивно, провівши користувацьке дослідження.
%     \end{itemize}
% \end{frame}